\documentclass[11pt,a4paper]{article}
% lesson-preamble.tex -- Shared preamble for all Phase 00 lessons
% Usage: % lesson-preamble.tex -- Shared preamble for all Phase 00 lessons
% Usage: % lesson-preamble.tex -- Shared preamble for all Phase 00 lessons
% Usage: \input{../lesson-preamble} (from subdomain directories)

% ---- Encoding and fonts ----
\usepackage[utf8]{inputenc}
\usepackage[T1]{fontenc}

% ---- Mathematics ----
\usepackage{amsmath,amssymb,amsthm}
\usepackage{mathtools}

% ---- Page geometry ----
\usepackage[margin=1in,headheight=14pt]{geometry}

% ---- Hyperlinks ----
\usepackage[colorlinks=true,linkcolor=red,citecolor=blue,urlcolor=blue]{hyperref}

% ---- Lists ----
\usepackage{enumitem}

% ---- Colors and boxes ----
\usepackage{xcolor}
\usepackage[theorems,skins,breakable]{tcolorbox}

% ---- Header/Footer ----
\usepackage{fancyhdr}
\renewcommand{\headrulewidth}{0.4pt}

% ---- Tables ----
\usepackage{booktabs}

% ---- Bibliography ----
\usepackage{natbib}
\bibliographystyle{plainnat}

% --------------------------------------------------------------------------
% Theorem environments
% --------------------------------------------------------------------------
\theoremstyle{plain}
\newtheorem{theorem}{Theorem}[section]
\newtheorem{lemma}[theorem]{Lemma}
\newtheorem{proposition}[theorem]{Proposition}
\newtheorem{corollary}[theorem]{Corollary}

\theoremstyle{definition}
\newtheorem{definition}[theorem]{Definition}
\newtheorem{example}[theorem]{Example}
\newtheorem{exercise}[theorem]{Exercise}
\newtheorem{assumption}[theorem]{Assumption}

\theoremstyle{remark}
\newtheorem{remark}[theorem]{Remark}

% --------------------------------------------------------------------------
% Colored boxes
% --------------------------------------------------------------------------
\newtcolorbox{keyresult}[1][]{%
  colback=green!5!white,
  colframe=green!50!black,
  fonttitle=\bfseries,
  title={Key Result},
  breakable,
  #1
}

\newtcolorbox{rlconnection}[1][]{%
  colback=blue!5!white,
  colframe=blue!50!black,
  fonttitle=\bfseries,
  title={RL/ML Connection},
  breakable,
  #1
}

\newtcolorbox{warning}[1][]{%
  colback=red!5!white,
  colframe=red!50!black,
  fonttitle=\bfseries,
  title={Warning},
  breakable,
  #1
}

% --------------------------------------------------------------------------
% Custom commands -- number sets
% --------------------------------------------------------------------------
\newcommand{\R}{\mathbb{R}}
\newcommand{\N}{\mathbb{N}}
\newcommand{\Q}{\mathbb{Q}}
\newcommand{\Z}{\mathbb{Z}}
\newcommand{\C}{\mathbb{C}}
\newcommand{\F}{\mathbb{F}}

% --------------------------------------------------------------------------
% Custom commands -- probability and expectation
% --------------------------------------------------------------------------
\newcommand{\E}{\mathbb{E}}
\newcommand{\Prob}{\mathbb{P}}
\newcommand{\Var}{\operatorname{Var}}
\newcommand{\Cov}{\operatorname{Cov}}
\newcommand{\indicator}{\mathbf{1}}

% --------------------------------------------------------------------------
% Custom commands -- convergence arrows
% --------------------------------------------------------------------------
\newcommand{\cas}{\xrightarrow{\text{a.s.}}}
\newcommand{\cprob}{\xrightarrow{\Prob}}
\newcommand{\clp}[1]{\xrightarrow{L^{#1}}}
\newcommand{\cdist}{\xrightarrow{d}}

% --------------------------------------------------------------------------
% Custom commands -- common operators
% --------------------------------------------------------------------------
\DeclareMathOperator{\dom}{dom}
\DeclareMathOperator{\epi}{epi}
\DeclareMathOperator{\conv}{conv}
\DeclareMathOperator{\relint}{relint}
\DeclareMathOperator{\aff}{aff}
\DeclareMathOperator{\cl}{cl}
\DeclareMathOperator{\interior}{int}
\DeclareMathOperator{\tr}{tr}
\DeclareMathOperator{\diag}{diag}
\DeclareMathOperator{\rank}{rank}
\DeclareMathOperator{\sgn}{sgn}
\DeclareMathOperator{\supp}{supp}
\DeclareMathOperator{\argmin}{arg\,min}
\DeclareMathOperator{\argmax}{arg\,max}
\DeclareMathOperator{\prox}{prox}
\DeclareMathOperator{\dist}{dist}
\DeclareMathOperator{\diam}{diam}
\DeclareMathOperator{\Span}{span}

% --------------------------------------------------------------------------
% Custom commands -- linear algebra operators
% --------------------------------------------------------------------------
\DeclareMathOperator{\nullity}{nullity}
\DeclareMathOperator{\range}{range}
\DeclareMathOperator{\nullspace}{null}
\DeclareMathOperator{\proj}{proj}

% --------------------------------------------------------------------------
% Custom commands -- measure theory
% --------------------------------------------------------------------------
\newcommand{\powerset}{\mathcal{P}}
\newcommand{\Borel}{\mathcal{B}}
 (from subdomain directories)

% ---- Encoding and fonts ----
\usepackage[utf8]{inputenc}
\usepackage[T1]{fontenc}

% ---- Mathematics ----
\usepackage{amsmath,amssymb,amsthm}
\usepackage{mathtools}

% ---- Page geometry ----
\usepackage[margin=1in,headheight=14pt]{geometry}

% ---- Hyperlinks ----
\usepackage[colorlinks=true,linkcolor=red,citecolor=blue,urlcolor=blue]{hyperref}

% ---- Lists ----
\usepackage{enumitem}

% ---- Colors and boxes ----
\usepackage{xcolor}
\usepackage[theorems,skins,breakable]{tcolorbox}

% ---- Header/Footer ----
\usepackage{fancyhdr}
\renewcommand{\headrulewidth}{0.4pt}

% ---- Tables ----
\usepackage{booktabs}

% ---- Bibliography ----
\usepackage{natbib}
\bibliographystyle{plainnat}

% --------------------------------------------------------------------------
% Theorem environments
% --------------------------------------------------------------------------
\theoremstyle{plain}
\newtheorem{theorem}{Theorem}[section]
\newtheorem{lemma}[theorem]{Lemma}
\newtheorem{proposition}[theorem]{Proposition}
\newtheorem{corollary}[theorem]{Corollary}

\theoremstyle{definition}
\newtheorem{definition}[theorem]{Definition}
\newtheorem{example}[theorem]{Example}
\newtheorem{exercise}[theorem]{Exercise}
\newtheorem{assumption}[theorem]{Assumption}

\theoremstyle{remark}
\newtheorem{remark}[theorem]{Remark}

% --------------------------------------------------------------------------
% Colored boxes
% --------------------------------------------------------------------------
\newtcolorbox{keyresult}[1][]{%
  colback=green!5!white,
  colframe=green!50!black,
  fonttitle=\bfseries,
  title={Key Result},
  breakable,
  #1
}

\newtcolorbox{rlconnection}[1][]{%
  colback=blue!5!white,
  colframe=blue!50!black,
  fonttitle=\bfseries,
  title={RL/ML Connection},
  breakable,
  #1
}

\newtcolorbox{warning}[1][]{%
  colback=red!5!white,
  colframe=red!50!black,
  fonttitle=\bfseries,
  title={Warning},
  breakable,
  #1
}

% --------------------------------------------------------------------------
% Custom commands -- number sets
% --------------------------------------------------------------------------
\newcommand{\R}{\mathbb{R}}
\newcommand{\N}{\mathbb{N}}
\newcommand{\Q}{\mathbb{Q}}
\newcommand{\Z}{\mathbb{Z}}
\newcommand{\C}{\mathbb{C}}
\newcommand{\F}{\mathbb{F}}

% --------------------------------------------------------------------------
% Custom commands -- probability and expectation
% --------------------------------------------------------------------------
\newcommand{\E}{\mathbb{E}}
\newcommand{\Prob}{\mathbb{P}}
\newcommand{\Var}{\operatorname{Var}}
\newcommand{\Cov}{\operatorname{Cov}}
\newcommand{\indicator}{\mathbf{1}}

% --------------------------------------------------------------------------
% Custom commands -- convergence arrows
% --------------------------------------------------------------------------
\newcommand{\cas}{\xrightarrow{\text{a.s.}}}
\newcommand{\cprob}{\xrightarrow{\Prob}}
\newcommand{\clp}[1]{\xrightarrow{L^{#1}}}
\newcommand{\cdist}{\xrightarrow{d}}

% --------------------------------------------------------------------------
% Custom commands -- common operators
% --------------------------------------------------------------------------
\DeclareMathOperator{\dom}{dom}
\DeclareMathOperator{\epi}{epi}
\DeclareMathOperator{\conv}{conv}
\DeclareMathOperator{\relint}{relint}
\DeclareMathOperator{\aff}{aff}
\DeclareMathOperator{\cl}{cl}
\DeclareMathOperator{\interior}{int}
\DeclareMathOperator{\tr}{tr}
\DeclareMathOperator{\diag}{diag}
\DeclareMathOperator{\rank}{rank}
\DeclareMathOperator{\sgn}{sgn}
\DeclareMathOperator{\supp}{supp}
\DeclareMathOperator{\argmin}{arg\,min}
\DeclareMathOperator{\argmax}{arg\,max}
\DeclareMathOperator{\prox}{prox}
\DeclareMathOperator{\dist}{dist}
\DeclareMathOperator{\diam}{diam}
\DeclareMathOperator{\Span}{span}

% --------------------------------------------------------------------------
% Custom commands -- linear algebra operators
% --------------------------------------------------------------------------
\DeclareMathOperator{\nullity}{nullity}
\DeclareMathOperator{\range}{range}
\DeclareMathOperator{\nullspace}{null}
\DeclareMathOperator{\proj}{proj}

% --------------------------------------------------------------------------
% Custom commands -- measure theory
% --------------------------------------------------------------------------
\newcommand{\powerset}{\mathcal{P}}
\newcommand{\Borel}{\mathcal{B}}
 (from subdomain directories)

% ---- Encoding and fonts ----
\usepackage[utf8]{inputenc}
\usepackage[T1]{fontenc}

% ---- Mathematics ----
\usepackage{amsmath,amssymb,amsthm}
\usepackage{mathtools}

% ---- Page geometry ----
\usepackage[margin=1in,headheight=14pt]{geometry}

% ---- Hyperlinks ----
\usepackage[colorlinks=true,linkcolor=red,citecolor=blue,urlcolor=blue]{hyperref}

% ---- Lists ----
\usepackage{enumitem}

% ---- Colors and boxes ----
\usepackage{xcolor}
\usepackage[theorems,skins,breakable]{tcolorbox}

% ---- Header/Footer ----
\usepackage{fancyhdr}
\renewcommand{\headrulewidth}{0.4pt}

% ---- Tables ----
\usepackage{booktabs}

% ---- Bibliography ----
\usepackage{natbib}
\bibliographystyle{plainnat}

% --------------------------------------------------------------------------
% Theorem environments
% --------------------------------------------------------------------------
\theoremstyle{plain}
\newtheorem{theorem}{Theorem}[section]
\newtheorem{lemma}[theorem]{Lemma}
\newtheorem{proposition}[theorem]{Proposition}
\newtheorem{corollary}[theorem]{Corollary}

\theoremstyle{definition}
\newtheorem{definition}[theorem]{Definition}
\newtheorem{example}[theorem]{Example}
\newtheorem{exercise}[theorem]{Exercise}
\newtheorem{assumption}[theorem]{Assumption}

\theoremstyle{remark}
\newtheorem{remark}[theorem]{Remark}

% --------------------------------------------------------------------------
% Colored boxes
% --------------------------------------------------------------------------
\newtcolorbox{keyresult}[1][]{%
  colback=green!5!white,
  colframe=green!50!black,
  fonttitle=\bfseries,
  title={Key Result},
  breakable,
  #1
}

\newtcolorbox{rlconnection}[1][]{%
  colback=blue!5!white,
  colframe=blue!50!black,
  fonttitle=\bfseries,
  title={RL/ML Connection},
  breakable,
  #1
}

\newtcolorbox{warning}[1][]{%
  colback=red!5!white,
  colframe=red!50!black,
  fonttitle=\bfseries,
  title={Warning},
  breakable,
  #1
}

% --------------------------------------------------------------------------
% Custom commands -- number sets
% --------------------------------------------------------------------------
\newcommand{\R}{\mathbb{R}}
\newcommand{\N}{\mathbb{N}}
\newcommand{\Q}{\mathbb{Q}}
\newcommand{\Z}{\mathbb{Z}}
\newcommand{\C}{\mathbb{C}}
\newcommand{\F}{\mathbb{F}}

% --------------------------------------------------------------------------
% Custom commands -- probability and expectation
% --------------------------------------------------------------------------
\newcommand{\E}{\mathbb{E}}
\newcommand{\Prob}{\mathbb{P}}
\newcommand{\Var}{\operatorname{Var}}
\newcommand{\Cov}{\operatorname{Cov}}
\newcommand{\indicator}{\mathbf{1}}

% --------------------------------------------------------------------------
% Custom commands -- convergence arrows
% --------------------------------------------------------------------------
\newcommand{\cas}{\xrightarrow{\text{a.s.}}}
\newcommand{\cprob}{\xrightarrow{\Prob}}
\newcommand{\clp}[1]{\xrightarrow{L^{#1}}}
\newcommand{\cdist}{\xrightarrow{d}}

% --------------------------------------------------------------------------
% Custom commands -- common operators
% --------------------------------------------------------------------------
\DeclareMathOperator{\dom}{dom}
\DeclareMathOperator{\epi}{epi}
\DeclareMathOperator{\conv}{conv}
\DeclareMathOperator{\relint}{relint}
\DeclareMathOperator{\aff}{aff}
\DeclareMathOperator{\cl}{cl}
\DeclareMathOperator{\interior}{int}
\DeclareMathOperator{\tr}{tr}
\DeclareMathOperator{\diag}{diag}
\DeclareMathOperator{\rank}{rank}
\DeclareMathOperator{\sgn}{sgn}
\DeclareMathOperator{\supp}{supp}
\DeclareMathOperator{\argmin}{arg\,min}
\DeclareMathOperator{\argmax}{arg\,max}
\DeclareMathOperator{\prox}{prox}
\DeclareMathOperator{\dist}{dist}
\DeclareMathOperator{\diam}{diam}
\DeclareMathOperator{\Span}{span}

% --------------------------------------------------------------------------
% Custom commands -- linear algebra operators
% --------------------------------------------------------------------------
\DeclareMathOperator{\nullity}{nullity}
\DeclareMathOperator{\range}{range}
\DeclareMathOperator{\nullspace}{null}
\DeclareMathOperator{\proj}{proj}

% --------------------------------------------------------------------------
% Custom commands -- measure theory
% --------------------------------------------------------------------------
\newcommand{\powerset}{\mathcal{P}}
\newcommand{\Borel}{\mathcal{B}}

\pagestyle{fancy}
\fancyhf{}
\fancyhead[L]{\textsc{Lesson 1: Real Numbers, Sequences, and Series}}
\fancyhead[R]{\thepage}
\title{Lesson 1: Real Numbers, Sequences, and Series}
\author{AI/ML Mathematics}
\date{}
\begin{document}
\maketitle
\tableofcontents
\newpage

%% ============================================================
%% SECTION 1: MOTIVATION AND OVERVIEW
%% ============================================================
\section{Motivation and Overview}
\label{sec:motivation}

Consider a Markov decision process with state space $S$ and discount factor
$\gamma \in (0,1)$. The Bellman optimality operator $T$ maps any bounded
value function $V$ to the function
\[
  (TV)(s) = \max_{a} \Bigl[ r(s,a)
    + \gamma \sum_{s'} p(s' \mid s,a)\, V(s') \Bigr].
\]
Value iteration generates the sequence $V_0, V_1 = TV_0, V_2 = TV_1, \ldots$
and one asserts that $V_n$ converges to the unique optimal value function
$V^*$. But \emph{why} does this limit exist? The proof relies on the Banach
fixed-point theorem, which itself rests on the completeness of $\R$: every
Cauchy sequence of real numbers converges. Without the completeness axiom,
the iterates could approach a limit that simply does not exist in our number
system---exactly the situation one encounters in $\Q$. The tools developed in
this lesson will let us formalize convergence, suprema, and the completeness
property that underwrites every convergence argument in reinforcement
learning and optimization.

\begin{rlconnection}[title={Why RL/ML needs the real number system}]
\begin{itemize}[nosep]
  \item \textbf{Value iteration.} Convergence of the sequence $V_n = T^n V_0$
    to $V^*$ requires $\R$ to be complete (every Cauchy sequence converges).
  \item \textbf{Gradient descent.} The statement ``$f(x_k)$ converges to
    $\inf_x f(x)$'' relies on the infimum existing, guaranteed by the
    completeness axiom.
  \item \textbf{Monte Carlo estimation.} The law of large numbers asserts
    $\frac{1}{n}\sum_{i=1}^n X_i \to \E[X]$; its proof uses the
    Bolzano--Weierstrass theorem and Cauchy completeness developed here.
  \item \textbf{Series in function approximation.} Neural network
    universality theorems express target functions as convergent series;
    the ratio and root tests verify absolute convergence of these
    expansions.
\end{itemize}
\end{rlconnection}

\paragraph{Prerequisites.}
Familiarity with high-school algebra (field operations, inequalities,
absolute values), basic set theory (subsets, unions, intersections), and
informal notions of natural numbers $\N$, integers $\Z$, and rationals $\Q$.
No prior analysis course is assumed.

\paragraph{Learning objectives.}
After completing this lesson, you will be able to:
\begin{itemize}[nosep]
  \item \textbf{State} the ordered-field axioms and the completeness axiom
    for $\R$, and \textbf{explain} why $\Q$ fails to be complete.
  \item \textbf{Define} supremum, infimum, convergence of sequences, and
    Cauchy sequences.
  \item \textbf{Prove} the Archimedean property, the monotone convergence
    theorem, and the Bolzano--Weierstrass theorem.
  \item \textbf{Compute} limits using the algebraic limit theorems and
    determine convergence of series via the ratio and root tests.
  \item \textbf{Derive} the relationship between $\limsup$, $\liminf$,
    and convergence of a bounded sequence.
\end{itemize}

%% ============================================================
%% SECTION 2: REFERENCES
%% ============================================================
\section{References}
\label{sec:references}

\begin{itemize}
  \item \citet[Chapters 1--3]{rudin1976} --- The canonical reference for
    the real number system, sequences, and series; concise and rigorous.
  \item \citet[Chapters 1--2]{abbott2015} --- An accessible introduction
    emphasizing the completeness axiom and its consequences.
  \item \citet[Chapter 0]{folland1999} --- A brief review of real-number
    foundations aimed at students proceeding to measure theory.
  \item \citet[Chapters 2--3]{bartle2011} --- Detailed treatment of
    sequences, series, and convergence with many worked examples.
  \item \citet[Chapter 1]{tao2016} --- A construction-oriented approach
    building $\R$ from the rationals via Cauchy sequences.
\end{itemize}

%% ============================================================
%% SECTION 3: THE REAL NUMBER SYSTEM
%% ============================================================
\section{The Real Number System}
\label{sec:real_numbers}

\begin{definition}[Ordered Field]
\label{def:ordered_field}
The real numbers $\R$ form an \emph{ordered field}: a set equipped with
operations $+$ and $\cdot$ and a total order $\le$ satisfying:
\begin{enumerate}[nosep,label=(\roman*)]
  \item \textbf{Field axioms.} $(\R, +, \cdot)$ is a field: addition and
    multiplication are commutative, associative, and distributive, with
    identity elements $0$ and $1$ respectively, additive inverses for all
    elements, and multiplicative inverses for all nonzero elements.
  \item \textbf{Order axioms.} The relation $\le$ is a total order
    compatible with the field operations: if $a \le b$ then
    $a + c \le b + c$ for all $c$; if $0 \le a$ and $0 \le b$ then
    $0 \le ab$.
\end{enumerate}
\end{definition}

\begin{remark}
\label{rem:Q_ordered_field}
The rationals $\Q$ also form an ordered field. What distinguishes $\R$
from $\Q$ is the completeness axiom below.
\end{remark}

\begin{example}
\label{ex:Q_incomplete}
The set $S = \{q \in \Q : q^2 < 2\}$ is nonempty (e.g., $1 \in S$) and
bounded above in $\Q$ (e.g., by $2$). However, $S$ has no supremum
\emph{in} $\Q$: any candidate rational upper bound can be improved by
moving closer to $\sqrt{2}$, which is irrational. This failure motivates
the completeness axiom.
\end{example}

\begin{definition}[Upper Bound, Supremum]
\label{def:supremum}
Let $S \subseteq \R$ be nonempty.
\begin{enumerate}[nosep,label=(\roman*)]
  \item An element $u \in \R$ is an \emph{upper bound} for $S$ if
    $s \le u$ for every $s \in S$. We say $S$ is \emph{bounded above}.
  \item An element $\alpha \in \R$ is the \emph{supremum} (or
    \emph{least upper bound}) of $S$, written $\alpha = \sup S$, if:
    \begin{enumerate}[nosep,label=(\alph*)]
      \item $\alpha$ is an upper bound for $S$, and
      \item if $u$ is any upper bound for $S$, then $\alpha \le u$.
    \end{enumerate}
\end{enumerate}
\end{definition}

\begin{definition}[Lower Bound, Infimum]
\label{def:infimum}
Dually, $\ell \in \R$ is a \emph{lower bound} for $S$ if $\ell \le s$ for
every $s \in S$. The \emph{infimum} $\inf S$ is the greatest lower bound:
$\inf S$ is a lower bound and $\ell \le \inf S$ for every lower bound
$\ell$.
\end{definition}

\begin{remark}
\label{rem:sup_epsilon}
The supremum $\alpha = \sup S$ is characterized by the
\emph{$\varepsilon$-criterion}: $\alpha$ is an upper bound for $S$, and
for every $\varepsilon > 0$ there exists $s \in S$ with
$s > \alpha - \varepsilon$. This characterization is used constantly in
analysis.
\end{remark}

\begin{keyresult}[title={Completeness Axiom}]
\begin{theorem}[Completeness Axiom / Least Upper Bound Property]
\label{thm:completeness}
Every nonempty subset of\/ $\R$ that is bounded above has a supremum
in~$\R$.
\end{theorem}
This is an \emph{axiom}: it is not proved from more basic principles but
is the defining property that distinguishes $\R$ from $\Q$. It
guarantees that limits, infima, and suprema exist whenever they should.
\end{keyresult}

\begin{corollary}[Greatest Lower Bound Property]
\label{cor:glb}
Every nonempty subset of\/ $\R$ that is bounded below has an infimum
in~$\R$.
\end{corollary}

\begin{proof}
Let $S \subseteq \R$ be nonempty and bounded below. Define
$-S = \{-s : s \in S\}$. Then $-S$ is nonempty and bounded above (if
$\ell$ is a lower bound for $S$, then $-\ell$ is an upper bound for $-S$).
By the completeness axiom, $\alpha = \sup(-S)$ exists. We claim
$\inf S = -\alpha$.

Since $\alpha$ is an upper bound for $-S$, we have $-s \le \alpha$ for
all $s \in S$, hence $s \ge -\alpha$, so $-\alpha$ is a lower bound for
$S$. If $\ell$ is any lower bound for $S$, then $-\ell$ is an upper
bound for $-S$, so $\alpha \le -\ell$, hence $\ell \le -\alpha$. Thus
$-\alpha$ is the greatest lower bound.
\end{proof}

\begin{rlconnection}[title={Suprema and infima in optimization}]
In optimization and RL, one routinely writes $\inf_{x \in X} f(x)$ or
$\sup_\pi V^\pi(s)$. The completeness axiom and the greatest lower bound
property guarantee these quantities exist as real numbers (when the sets
are bounded), rather than as formal symbols with no concrete value.
Without completeness, the optimal value function $V^*(s) = \sup_\pi
V^\pi(s)$ might fail to exist, and no algorithm could converge to it.
\end{rlconnection}

\begin{theorem}[Archimedean Property]
\label{thm:archimedean}
For every $x \in \R$, there exists $n \in \N$ with $n > x$.
Equivalently, for every $\varepsilon > 0$ and every $M > 0$, there exists
$n \in \N$ with $n\varepsilon > M$.
\end{theorem}

\begin{proof}
Suppose for contradiction that $\N$ is bounded above in $\R$. By the
completeness axiom, $\alpha = \sup \N$ exists. Then $\alpha - 1$ is not
an upper bound for $\N$ (since $\alpha$ is the \emph{least} upper
bound), so there exists $m \in \N$ with $m > \alpha - 1$. But then
$m + 1 > \alpha$ with $m + 1 \in \N$, contradicting the fact that
$\alpha$ is an upper bound for $\N$.

The second formulation follows by applying the first to
$x = M/\varepsilon$.
\end{proof}

\begin{example}
\label{ex:archimedean_step_size}
In gradient descent with step size $\eta > 0$, one often needs
$n$ iterations so that $n\eta$ exceeds a given threshold $M$
(e.g., to traverse a distance $M$ in parameter space). The Archimedean
property guarantees such $n$ exists for any positive step size, no matter
how small.
\end{example}

\begin{theorem}[Density of Rationals]
\label{thm:density_Q}
For every $a, b \in \R$ with $a < b$, there exists $q \in \Q$ with
$a < q < b$.
\end{theorem}

\begin{proof}
By the Archimedean property, choose $n \in \N$ with $n(b - a) > 1$,
i.e., $nb - na > 1$. Again by the Archimedean property, the sets
$\{k \in \Z : k > na\}$ and $\{k \in \Z : k < na\}$ are both nonempty,
so we may define $m = \min\{k \in \Z : k > na\}$ (this minimum exists
because the positive integers exceeding $na$ form a nonempty subset of
$\N$, which has a least element by the well-ordering principle). Then
$m > na$ and $m - 1 \le na$, so $m \le na + 1 < nb$ (since
$nb - na > 1$). Setting $q = m/n$ gives $a < q < b$ with $q \in \Q$.
\end{proof}

\begin{theorem}[Nested Intervals Property]
\label{thm:nested_intervals}
Let $I_n = [a_n, b_n]$ be a sequence of closed bounded intervals with
$I_1 \supseteq I_2 \supseteq I_3 \supseteq \cdots$. Then
$\bigcap_{n=1}^\infty I_n \neq \emptyset$.
\end{theorem}

\begin{proof}
The nesting condition gives $a_1 \le a_2 \le \cdots$ and
$b_1 \ge b_2 \ge \cdots$ with $a_n \le b_n$ for all $n$. In particular,
$a_n \le b_1$ for all $n$, so $\{a_n : n \in \N\}$ is bounded above.
By the completeness axiom, $\alpha = \sup\{a_n : n \in \N\}$ exists.

For every $n$, we have $a_n \le \alpha$ (since $\alpha$ is an upper
bound) and $\alpha \le b_n$ (since $b_n$ is an upper bound for
$\{a_k\}$: for $k \le n$, $a_k \le a_n \le b_n$; for $k > n$,
$a_k \le b_k \le b_n$). Therefore $\alpha \in [a_n, b_n]$ for every
$n$, giving $\alpha \in \bigcap_{n=1}^\infty I_n$.
\end{proof}

\begin{warning}[title={Closure and boundedness are both essential}]
The nested intervals conclusion fails for open intervals:
$\bigcap_{n=1}^\infty (0, 1/n) = \emptyset$.
It also fails for unbounded closed sets:
$\bigcap_{n=1}^\infty [n, \infty) = \emptyset$.
Both closure \emph{and} boundedness are required.
\end{warning}

\begin{example}
\label{ex:bisection}
The bisection method for root-finding constructs nested intervals
$[a_n, b_n]$ with $f(a_n) < 0 < f(b_n)$ and $b_n - a_n = (b_0 - a_0)/2^n$.
The nested intervals property guarantees a point in $\bigcap_n [a_n, b_n]$,
and continuity of $f$ forces $f$ to vanish there.
\end{example}

With the real number system and its completeness properties established,
we now turn to the central objects of analysis: sequences and their
convergence behavior.

%% ============================================================
%% SECTION 4: SEQUENCES
%% ============================================================
\section{Sequences}
\label{sec:sequences}

\begin{definition}[Convergence]
\label{def:convergence}
A sequence $(a_n)_{n=1}^\infty$ in $\R$ \emph{converges} to $L \in \R$,
written $\lim_{n \to \infty} a_n = L$ or $a_n \to L$, if for every
$\varepsilon > 0$ there exists $N \in \N$ such that $|a_n - L| <
\varepsilon$ for all $n \ge N$. A sequence that does not converge is
\emph{divergent}.
\end{definition}

\begin{example}
\label{ex:seq_1_over_n}
The sequence $a_n = 1/n$ converges to $0$. Given $\varepsilon > 0$,
choose $N \in \N$ with $N > 1/\varepsilon$ (Archimedean property). For
$n \ge N$: $|a_n - 0| = 1/n \le 1/N < \varepsilon$.
\end{example}

\begin{theorem}[Uniqueness of Limits]
\label{thm:unique_limit}
If $a_n \to L$ and $a_n \to M$, then $L = M$.
\end{theorem}

\begin{proof}
Let $\varepsilon > 0$. There exist $N_1, N_2 \in \N$ such that
$|a_n - L| < \varepsilon/2$ for $n \ge N_1$ and $|a_n - M| <
\varepsilon/2$ for $n \ge N_2$. For $n \ge \max(N_1, N_2)$:
\[
  |L - M| \le |L - a_n| + |a_n - M|
  < \frac{\varepsilon}{2} + \frac{\varepsilon}{2} = \varepsilon.
\]
Since $|L - M| < \varepsilon$ for every $\varepsilon > 0$, we conclude
$L = M$.
\end{proof}

\begin{theorem}[Algebraic Limit Theorems]
\label{thm:alg_limit}
Suppose $a_n \to a$ and $b_n \to b$. Then:
\begin{enumerate}[nosep,label=(\roman*)]
  \item $a_n + b_n \to a + b$.
  \item $a_n \cdot b_n \to a \cdot b$.
  \item If $b \neq 0$, then $a_n / b_n \to a / b$ (for $n$ sufficiently
    large that $b_n \neq 0$).
\end{enumerate}
\end{theorem}

\begin{proof}
We prove (ii); parts (i) and (iii) follow by similar arguments.

Since $(a_n)$ converges, it is bounded: there exists $M > 0$ with
$|a_n| \le M$ for all $n$. Let $\varepsilon > 0$. Choose $N_1$ so that
$|a_n - a| < \varepsilon/(2(|b| + 1))$ for $n \ge N_1$, and $N_2$ so
that $|b_n - b| < \varepsilon/(2M)$ for $n \ge N_2$. For
$n \ge \max(N_1, N_2)$:
\begin{align*}
  |a_n b_n - ab|
  &= |a_n b_n - a_n b + a_n b - ab| \\
  &\le |a_n||b_n - b| + |b||a_n - a| \\
  &\le M \cdot \frac{\varepsilon}{2M}
       + |b| \cdot \frac{\varepsilon}{2(|b|+1)} \\
  &< \frac{\varepsilon}{2} + \frac{\varepsilon}{2} = \varepsilon.
    \qedhere
\end{align*}
\end{proof}

\begin{example}
\label{ex:alg_limit}
Consider the sequence $a_n = (3n^2 + 1)/(n^2 + 5n)$. Dividing numerator
and denominator by $n^2$:
\[
  a_n = \frac{3 + 1/n^2}{1 + 5/n}.
\]
Since $1/n^2 \to 0$ and $5/n \to 0$, the algebraic limit theorems give
$a_n \to 3/1 = 3$.
\end{example}

\begin{keyresult}[title={Monotone Convergence Theorem}]
\begin{theorem}[Monotone Convergence Theorem for Sequences]
\label{thm:mct_seq}
Every bounded monotone sequence converges. Specifically:
\begin{enumerate}[nosep,label=(\roman*)]
  \item If $(a_n)$ is increasing and bounded above, then
    $a_n \to \sup\{a_n : n \in \N\}$.
  \item If $(a_n)$ is decreasing and bounded below, then
    $a_n \to \inf\{a_n : n \in \N\}$.
\end{enumerate}
\end{theorem}
\end{keyresult}

\begin{proof}
We prove (i). Let $\alpha = \sup\{a_n : n \in \N\}$, which exists by the
completeness axiom. Let $\varepsilon > 0$. Since $\alpha - \varepsilon$
is not an upper bound for $\{a_n\}$, there exists $N$ with
$a_N > \alpha - \varepsilon$. Since $(a_n)$ is increasing, for
$n \ge N$:
\[
  \alpha - \varepsilon < a_N \le a_n \le \alpha,
\]
so $|a_n - \alpha| < \varepsilon$. Part (ii) follows by applying (i) to
$(-a_n)$.
\end{proof}

\begin{rlconnection}[title={Monotone convergence in policy iteration}]
Policy iteration generates a sequence of value functions $V^{\pi_0},
V^{\pi_1}, V^{\pi_2}, \ldots$ that is monotonically non-decreasing
(each policy improvement step satisfies $V^{\pi_{k+1}}(s) \ge
V^{\pi_k}(s)$ for all $s$) and bounded above by $V^*(s) \le
R_{\max}/(1-\gamma)$. The monotone convergence theorem guarantees
convergence of this sequence.
\end{rlconnection}

\begin{theorem}[Bolzano--Weierstrass]
\label{thm:bolzano_weierstrass}
Every bounded sequence in $\R$ has a convergent subsequence.
\end{theorem}

\begin{proof}
Let $(a_n)$ be bounded. We construct a monotone subsequence, which
converges by Theorem~\ref{thm:mct_seq}.

Call an index $n$ a \emph{peak} if $a_n \ge a_m$ for all $m > n$.

\textbf{Case 1:} There are infinitely many peaks $n_1 < n_2 < n_3 <
\cdots$. Then $a_{n_1} \ge a_{n_2} \ge a_{n_3} \ge \cdots$ is a
decreasing subsequence. It is bounded (since $(a_n)$ is bounded), hence
convergent by Theorem~\ref{thm:mct_seq}.

\textbf{Case 2:} There are only finitely many peaks. Let $N$ be larger
than all peak indices. Then no index $n \ge N$ is a peak, meaning for
every $n \ge N$ there exists $m > n$ with $a_m > a_n$. Starting with
$n_1 = N$, inductively choose $n_{k+1} > n_k$ with
$a_{n_{k+1}} > a_{n_k}$. This gives a strictly increasing subsequence.
It is bounded above (since $(a_n)$ is bounded), hence convergent by
Theorem~\ref{thm:mct_seq}.
\end{proof}

\begin{example}
\label{ex:subsequence}
The sequence $a_n = (-1)^n(1 + 1/n)$ is bounded but divergent. However,
the subsequence $a_{2k} = 1 + 1/(2k) \to 1$ converges, as guaranteed by
Bolzano--Weierstrass.
\end{example}

\begin{definition}[Cauchy Sequence]
\label{def:cauchy}
A sequence $(a_n)$ in $\R$ is \emph{Cauchy} if for every
$\varepsilon > 0$ there exists $N \in \N$ such that
$|a_n - a_m| < \varepsilon$ for all $n, m \ge N$.
\end{definition}

\begin{theorem}[Completeness of $\R$: Cauchy Criterion]
\label{thm:cauchy_complete}
A sequence in $\R$ converges if and only if it is Cauchy.
\end{theorem}

\begin{proof}
$(\Rightarrow)$ Suppose $a_n \to L$. Given $\varepsilon > 0$, choose $N$
so that $|a_n - L| < \varepsilon/2$ for $n \ge N$. For $n, m \ge N$:
\[
  |a_n - a_m| \le |a_n - L| + |L - a_m| < \varepsilon.
\]

$(\Leftarrow)$ Suppose $(a_n)$ is Cauchy. First, $(a_n)$ is bounded:
choose $N$ so that $|a_n - a_m| < 1$ for $n, m \ge N$; then
$|a_n| \le \max(|a_1|, \ldots, |a_{N-1}|, |a_N| + 1)$ for all $n$.

By Bolzano--Weierstrass (Theorem~\ref{thm:bolzano_weierstrass}),
$(a_n)$ has a convergent subsequence $a_{n_k} \to L$. We show
$a_n \to L$. Given $\varepsilon > 0$, choose $N_1$ so that
$|a_n - a_m| < \varepsilon/2$ for $n, m \ge N_1$, and choose $K$ so
that $|a_{n_k} - L| < \varepsilon/2$ for $k \ge K$. Pick $k \ge K$
with $n_k \ge N_1$. Then for $n \ge N_1$:
\[
  |a_n - L| \le |a_n - a_{n_k}| + |a_{n_k} - L|
  < \frac{\varepsilon}{2} + \frac{\varepsilon}{2} = \varepsilon.
  \qedhere
\]
\end{proof}

\begin{rlconnection}[title={Cauchy sequences in iterative algorithms}]
The Cauchy criterion is indispensable in RL and optimization because it
certifies convergence \emph{without knowing the limit}. In practice, we
often cannot compute $V^*$ or $\theta^*$ directly. Instead, we show
successive iterates $\|V_{n+1} - V_n\|$ shrink geometrically (e.g., by a
factor $\gamma$), verifying the Cauchy condition, and conclude that the
limit exists. This is precisely the logic behind contraction-mapping
proofs of value iteration convergence.
\end{rlconnection}

\begin{definition}[$\limsup$ and $\liminf$]
\label{def:limsup_liminf}
For a bounded sequence $(a_n)$, define:
\[
  \limsup_{n \to \infty} a_n
  = \lim_{n \to \infty}\Bigl(\sup_{k \ge n} a_k\Bigr), \qquad
  \liminf_{n \to \infty} a_n
  = \lim_{n \to \infty}\Bigl(\inf_{k \ge n} a_k\Bigr).
\]
The sequence $\sup_{k \ge n} a_k$ is decreasing and bounded below
(by $\inf a_k$), and $\inf_{k \ge n} a_k$ is increasing and bounded
above (by $\sup a_k$), so both limits exist by
Theorem~\ref{thm:mct_seq}.
\end{definition}

\begin{example}
\label{ex:limsup_liminf}
For $a_n = (-1)^n + 1/n$, the odd-indexed terms approach $-1$ from above
and the even-indexed terms approach $1$ from above. We have
$\sup_{k \ge n} a_k \to 1$ and $\inf_{k \ge n} a_k \to -1$, so
$\limsup a_n = 1$ and $\liminf a_n = -1$. Since these differ, the
sequence diverges.
\end{example}

\begin{proposition}[Properties of $\limsup$ and $\liminf$]
\label{prop:limsup_props}
Let $(a_n)$ be a bounded sequence. Then:
\begin{enumerate}[nosep,label=(\roman*)]
  \item $\liminf_{n\to\infty} a_n \le \limsup_{n\to\infty} a_n$.
  \item $(a_n)$ converges if and only if
    $\liminf_{n\to\infty} a_n = \limsup_{n\to\infty} a_n$, and in that
    case the common value equals $\lim a_n$.
  \item $\limsup_{n\to\infty} a_n = L$ if and only if: for every
    $\varepsilon > 0$, (a)~$a_n < L + \varepsilon$ for all sufficiently
    large $n$, and (b)~$a_n > L - \varepsilon$ for infinitely many $n$.
\end{enumerate}
\end{proposition}

\begin{proof}
(i) For each $n$, $\inf_{k \ge n} a_k \le \sup_{k \ge n} a_k$. Taking
limits preserves the inequality.

(ii) $(\Rightarrow)$ Suppose $a_n \to L$. Given $\varepsilon > 0$, for
large $n$ we have $L - \varepsilon < a_k < L + \varepsilon$ for all
$k \ge n$, so
$L - \varepsilon \le \inf_{k \ge n} a_k \le \sup_{k \ge n} a_k \le
L + \varepsilon$.
Taking $n \to \infty$ gives
$L - \varepsilon \le \liminf a_n \le \limsup a_n \le L + \varepsilon$
for all $\varepsilon > 0$, so both equal $L$.

$(\Leftarrow)$ Suppose $\liminf a_n = \limsup a_n = L$. For any
$\varepsilon > 0$, choose $N$ so that
$|\sup_{k \ge n} a_k - L| < \varepsilon$ and
$|\inf_{k \ge n} a_k - L| < \varepsilon$ for $n \ge N$. Then for
$n \ge N$:
\[
  L - \varepsilon < \inf_{k \ge n} a_k \le a_n
  \le \sup_{k \ge n} a_k < L + \varepsilon,
\]
so $|a_n - L| < \varepsilon$.

(iii) Let $s_n = \sup_{k \ge n} a_k$ and $L = \lim s_n$.

For (a): given $\varepsilon > 0$, choose $N$ with $s_N < L + \varepsilon$.
Then $a_n \le s_N < L + \varepsilon$ for all $n \ge N$.

For (b): given $\varepsilon > 0$, for every $n$ we have $s_n \ge L$
(since $s_n$ decreases to $L$), so there exists $k \ge n$ with
$a_k > s_n - \varepsilon \ge L - \varepsilon$. Since $n$ is arbitrary,
infinitely many such $k$ exist.

Conversely, suppose (a) and (b) hold for some value $L'$. From (a),
$s_n \le L' + \varepsilon$ for large $n$, so
$\limsup a_n \le L' + \varepsilon$. From (b), for each $n$ there exists
$k \ge n$ with $a_k > L' - \varepsilon$, so
$s_n \ge L' - \varepsilon$, giving
$\limsup a_n \ge L' - \varepsilon$. Since $\varepsilon$ is arbitrary,
$\limsup a_n = L'$.
\end{proof}

\begin{warning}[title={$\limsup$ is not $\sup$}]
A common mistake is to confuse $\limsup_{n\to\infty} a_n$ with
$\sup_n a_n$. The supremum $\sup_n a_n$ is the largest value attained
(or approached) by the entire sequence, while $\limsup a_n$ captures
the largest \emph{cluster point}---the largest value that subsequences
approach. For instance, the sequence $a_1 = 100, a_n = (-1)^n$ for
$n \ge 2$ has $\sup a_n = 100$ but $\limsup a_n = 1$.
\end{warning}

With the theory of sequences complete---including criteria for
convergence, extraction of convergent subsequences, and the $\limsup$ and
$\liminf$ tools---we are ready to study infinite series.

%% ============================================================
%% SECTION 5: SERIES
%% ============================================================
\section{Series}
\label{sec:series}

\begin{definition}[Convergence of Series and Absolute Convergence]
\label{def:series}
Let $(a_n)_{n=1}^\infty$ be a sequence. The \emph{series}
$\sum_{n=1}^\infty a_n$ \emph{converges} if the sequence of partial
sums $S_N = \sum_{n=1}^N a_n$ converges. The series \emph{converges
absolutely} if $\sum_{n=1}^\infty |a_n|$ converges.
\end{definition}

\begin{example}
\label{ex:geometric_series}
The geometric series $\sum_{n=0}^\infty r^n$ converges for $|r| < 1$.
Indeed, the partial sum $S_N = (1 - r^{N+1})/(1 - r)$ satisfies
$S_N \to 1/(1 - r)$, since $|r|^{N+1} \to 0$ by the Archimedean
property and algebraic limit theorems.
\end{example}

\begin{theorem}[Comparison Test]
\label{thm:comparison}
If $0 \le a_n \le b_n$ for all $n$ and $\sum b_n$ converges, then
$\sum a_n$ converges.
\end{theorem}

\begin{proof}
The partial sums $S_N = \sum_{n=1}^N a_n$ form an increasing sequence
(since $a_n \ge 0$) bounded above by $\sum_{n=1}^\infty b_n < \infty$.
By the monotone convergence theorem (Theorem~\ref{thm:mct_seq}),
$(S_N)$ converges.
\end{proof}

\begin{theorem}[Absolute Convergence Implies Convergence]
\label{thm:abs_conv}
If $\sum_{n=1}^\infty |a_n|$ converges, then $\sum_{n=1}^\infty a_n$
converges.
\end{theorem}

\begin{proof}
Absolute convergence means the partial sums of $|a_n|$ form a Cauchy
sequence. For $m > n$:
\[
  \Bigl|\sum_{k=n+1}^{m} a_k\Bigr| \le \sum_{k=n+1}^{m} |a_k|,
\]
so the partial sums of $a_n$ are also Cauchy, hence convergent by
Theorem~\ref{thm:cauchy_complete}.
\end{proof}

\begin{example}
\label{ex:abs_conv}
The series $\sum_{n=1}^\infty (-1)^n / n^2$ converges absolutely, since
$\sum 1/n^2$ converges (this can be shown by comparison with
$\sum 1/(n(n-1))$, which telescopes to $1$). By
Theorem~\ref{thm:abs_conv}, the original alternating series converges.
\end{example}

\begin{keyresult}[title={Ratio and Root Tests}]
\begin{theorem}[Ratio Test]
\label{thm:ratio}
Let $(a_n)$ be a sequence with $a_n \neq 0$ for all $n$.
\begin{enumerate}[nosep,label=(\roman*)]
  \item If $\limsup_{n\to\infty} |a_{n+1}/a_n| < 1$, then $\sum a_n$
    converges absolutely.
  \item If $\liminf_{n\to\infty} |a_{n+1}/a_n| > 1$, then $\sum a_n$
    diverges.
\end{enumerate}
\end{theorem}
\begin{theorem}[Root Test]
\label{thm:root}
Let $\alpha = \limsup_{n\to\infty} |a_n|^{1/n}$.
\begin{enumerate}[nosep,label=(\roman*)]
  \item If $\alpha < 1$, then $\sum a_n$ converges absolutely.
  \item If $\alpha > 1$, then $\sum a_n$ diverges.
\end{enumerate}
\end{theorem}
\end{keyresult}

\begin{proof}[Proof of the Ratio Test]
Suppose $\limsup |a_{n+1}/a_n| = L < 1$. Choose $r$ with $L < r < 1$.
By Proposition~\ref{prop:limsup_props}(iii), there exists $N$ such that
$|a_{n+1}/a_n| < r$ for all $n \ge N$. Then $|a_n| \le |a_N| r^{n-N}$
for $n \ge N$, and $\sum |a_N| r^{n-N}$ converges (geometric series
with ratio $r < 1$). By the comparison test, $\sum |a_n|$ converges.

If $\liminf |a_{n+1}/a_n| > 1$, then $|a_{n+1}| > |a_n|$ for all
sufficiently large $n$, so $a_n \not\to 0$ and the series diverges
(since convergence requires $a_n \to 0$).
\end{proof}

\begin{proof}[Proof of the Root Test]
If $\alpha < 1$, choose $r$ with $\alpha < r < 1$. By
Proposition~\ref{prop:limsup_props}(iii), $|a_n|^{1/n} < r$ for all
sufficiently large $n$, so $|a_n| < r^n$. The comparison test with the
geometric series $\sum r^n$ gives absolute convergence.

If $\alpha > 1$, then $|a_n|^{1/n} > 1$ for infinitely many $n$
(Proposition~\ref{prop:limsup_props}(iii)), so $|a_n| > 1$ infinitely
often, meaning $a_n \not\to 0$ and the series diverges.
\end{proof}

\begin{example}
\label{ex:ratio_test}
Consider the series $\sum_{n=0}^\infty x^n / n!$ (the exponential
series). The ratio test gives
$|a_{n+1}/a_n| = |x|/(n+1) \to 0 < 1$,
so the series converges absolutely for every $x \in \R$. This defines
$e^x$ as a convergent power series.
\end{example}

\begin{rlconnection}[title={Series convergence in function approximation}]
In supervised learning, kernel methods express predictions as
$f(x) = \sum_{n=1}^\infty \alpha_n \phi_n(x)$ for basis functions
$\phi_n$. Convergence of this series (and control of the approximation
error from truncating to finitely many terms) is established via the
comparison, ratio, or root tests applied to the coefficient sequence
$(\alpha_n)$.
\end{rlconnection}

\begin{remark}[Conditional Convergence]
\label{rem:conditional}
A series that converges but does not converge absolutely is called
\emph{conditionally convergent}. The alternating harmonic series
$\sum (-1)^{n+1}/n$ is the classic example. By the Riemann
rearrangement theorem, any conditionally convergent series can be
rearranged to converge to any prescribed real value, or to diverge.
This pathology motivates the focus on absolute convergence in
applications. In RL and numerical analysis, one almost always works
with absolutely convergent series to ensure robustness under reordering
of summation.
\end{remark}

%% ============================================================
%% SECTION 6: EXERCISES
%% ============================================================
\section{Exercises}
\label{sec:exercises}

\begin{exercise}[Supremum computation]
\textnormal{(\(*\))} \\
Let $S = \{1 - 1/n : n \in \N\}$. Prove that $\sup S = 1$ directly from
the definition of supremum (Definition~\ref{def:supremum}).
\end{exercise}

\begin{exercise}[Convergence from the definition]
\textnormal{(\(*\))} \\
Prove directly from Definition~\ref{def:convergence} that the sequence
$a_n = (2n+1)/(n+3)$ converges to $2$.
\end{exercise}

\begin{exercise}[Monotone convergence application]
\textnormal{(\(**\))} \\
Define the sequence $a_1 = 1$, $a_{n+1} = \sqrt{2 + a_n}$. Prove that
$(a_n)$ is increasing and bounded above, and find its limit.

\textit{Hint:} Show by induction that $a_n < 2$ for all $n$. For the
limit, if $a_n \to L$ then $L = \sqrt{2 + L}$.
\end{exercise}

\begin{exercise}[Cauchy criterion in RL]
\textnormal{(\(**\))} \\
Let $T\colon \R \to \R$ be a contraction mapping with factor
$\gamma \in (0,1)$: $|Tx - Ty| \le \gamma|x - y|$ for all $x,y \in \R$.
Define $x_0 \in \R$ arbitrarily and $x_{n+1} = Tx_n$. Prove that
$(x_n)$ is a Cauchy sequence, and conclude that it converges.

\textit{Hint:} Show $|x_{n+1} - x_n| \le \gamma^n |x_1 - x_0|$, then
use the geometric series formula for $|x_n - x_m|$.
\end{exercise}

\begin{exercise}[Series convergence tests]
\textnormal{(\(*\))} \\
Determine whether each series converges or diverges, and state which test
you use:
\begin{enumerate}[nosep,label=(\alph*)]
  \item $\displaystyle \sum_{n=1}^\infty \frac{3^n}{n!}$
  \item $\displaystyle \sum_{n=1}^\infty \frac{n}{2^n}$
  \item $\displaystyle \sum_{n=1}^\infty \frac{1}{n^{1/2}}$
\end{enumerate}
\end{exercise}

\begin{exercise}[Limsup and subsequential limits]
\textnormal{(\(***\))} \\
Let $(a_n)$ be a bounded sequence. Prove that there exists a subsequence
$(a_{n_k})$ with $a_{n_k} \to \limsup_{n\to\infty} a_n$.

\textit{Hint:} Use the characterization in
Proposition~\ref{prop:limsup_props}(iii) to inductively select $n_k$
with $|a_{n_k} - L| < 1/k$ where $L = \limsup a_n$.
\end{exercise}

%% ============================================================
%% SECTION 7: SUMMARY
%% ============================================================
\section{Summary}
\label{sec:summary}

\begin{itemize}
  \item The real numbers $\R$ form an ordered field satisfying the
    \emph{completeness axiom}: every nonempty subset bounded above has a
    supremum. This axiom is what distinguishes $\R$ from $\Q$.

  \item The completeness axiom implies the \emph{greatest lower bound
    property}: every nonempty subset bounded below has an infimum
    (Corollary~\ref{cor:glb}).

  \item The \emph{Archimedean property} ($\N$ is unbounded) and the
    \emph{density of rationals} are consequences of completeness.

  \item A sequence $(a_n)$ converges to $L$ if the terms eventually stay
    within any prescribed distance $\varepsilon$ of $L$; limits are
    unique and respect algebraic operations.

  \item The \emph{monotone convergence theorem} states that every bounded
    monotone sequence converges; this directly underpins policy iteration
    convergence in RL.

  \item The \emph{Bolzano--Weierstrass theorem} guarantees that every
    bounded sequence has a convergent subsequence, a cornerstone for
    compactness arguments.

  \item A sequence converges if and only if it is \emph{Cauchy}
    (completeness of $\R$). The Cauchy criterion enables convergence
    proofs for iterative algorithms without knowing the limit.

  \item The $\limsup$ and $\liminf$ of a bounded sequence always exist;
    the sequence converges exactly when these two values coincide.

  \item For series, the \emph{comparison test}, \emph{ratio test}, and
    \emph{root test} provide practical tools for establishing absolute
    convergence. Absolute convergence implies convergence and guards
    against rearrangement pathologies.
\end{itemize}

%% ============================================================
%% SECTION 8: FURTHER READING
%% ============================================================
\section{Further Reading}
\label{sec:further_reading}

\begin{itemize}
  \item \textbf{Lesson 2} (Continuity, Differentiation, and Integration)
    builds on the convergence theory developed here. The intermediate
    value theorem, mean value theorem, and Riemann integration all
    rely on completeness and the Bolzano--Weierstrass theorem.

  \item \textbf{Lesson 3} (Topology of the Real Line and Metric Spaces)
    generalizes convergence to abstract metric spaces. Compactness in
    metric spaces is the natural extension of the Bolzano--Weierstrass
    property.

  \item \textbf{Lesson 4} (Sequences and Series of Functions) treats
    pointwise and uniform convergence, power series, and the Weierstrass
    M-test---a direct generalization of the comparison test to function
    series.

  \item For a deeper treatment of the construction of $\R$ from $\Q$
    via Dedekind cuts or Cauchy sequences, see \citet[Chapter 1]{rudin1976}
    or \citet[Chapters 1--5]{tao2016}.

  \item Students proceeding to Phase 01 (measure theory) will find that
    the monotone convergence theorem and Cauchy completeness generalize
    to Lebesgue integration, where the monotone convergence theorem for
    integrals and the completeness of $L^p$ spaces play analogous roles.

  \item \citet{pugh2015} provides an enriched treatment with many
    challenging exercises and a geometric perspective on real analysis
    that complements the more algebraic approach of \citet{rudin1976}.
\end{itemize}

\bibliography{real-analysis-refs}

\end{document}
